\documentclass[12pt]{article}
\usepackage[margin=1in]{geometry}
\usepackage{amsmath,amsthm,amssymb,scrextend}
\usepackage{array}
\setlength{\parindent}{0pt}
\begin{document}
\newcommand{\ds}{\displaystyle}

\begin{align*}
\tan\left(\frac{\pi}{4}-\frac{x}{2}\right)
&=\frac{1-\tan\frac{x}{2}}{1+\tan\frac{x}{2}}\\
&=\frac{1-\sin x}{\cos x}\\
&=\frac1{\cos x}-\frac{\sin x}{\cos x}\\
&=\sec x-\tan x
\end{align*}

\bigskip\hrule\bigskip

\textbf{(i)} Using $(*)$ with $\ds x=\frac{\pi}{4}$,

$$\tan\frac{\pi}{8}=\sec\frac{\pi}{4}-\tan\frac{\pi}{4}=\sqrt2-1$$

\begin{align*}
\tan\frac{11\pi}{24}
&=\tan\left(\frac{\pi}{3}+\frac{\pi}{8}\right)\\
&=\frac{\sqrt3+\sqrt2-1}{1-\sqrt3(\sqrt2-1)}\\
&=\frac{\sqrt3+\sqrt2-1}{\sqrt3-\sqrt6+1}
\end{align*}

\bigskip\hrule\bigskip

\textbf{(ii)}

$$(\sqrt3-\sqrt6+1)(2+\sqrt2+\sqrt3+\sqrt6)=\sqrt3+\sqrt2-1$$


$$\therefore \frac{\sqrt3+\sqrt2-1}{\sqrt3-\sqrt6+1}=2+\sqrt2+\sqrt3+\sqrt6$$

\bigskip\hrule\bigskip

\textbf{(iii)}

Using $(*)$ with $\ds x=\frac{11\pi}{24}$,

$$\tan\frac\pi{48}=\sec\frac{11\pi}{24}-\tan\frac{11\pi}{24}$$

\begin{align*}
\sec\frac{11\pi}{24}
&=\sqrt{1+\tan^2\frac{11\pi}{24}}\\
&=\sqrt{1+(2+\sqrt2+\sqrt3+\sqrt6)^2}\\
&=\sqrt{16+10\sqrt2+8\sqrt3+6\sqrt6}
\end{align*}


From part \textbf{(ii)},  $\ds \tan\frac{11\pi}{24}=2+\sqrt2+\sqrt3+\sqrt6$


$$\therefore \tan\frac\pi{48}
=\sqrt{16+10\sqrt2+8\sqrt3+6\sqrt6}-2-\sqrt2-\sqrt3-\sqrt6$$

\bigskip
\hrule
\begin{center}
    This solution was submitted by \textbf{Jon}\\[4pt]
    Published and refined by \textbf{Jon}\\[8pt]
    \textit{Last updated: July 2, 2026}
\end{center}

\end{document}